\documentclass[11pt, aspectratio=169]{beamer}
\usepackage[style=authoryear-icomp, backend=biber]{biblatex}
\usetheme{jcar}

%\includeonlyframes{a} % Para compilar solo en la que se está trabajando

\newcommand\blfootnotee[1]{%
  \begingroup
  \renewcommand\thefootnote{}\footnote{#1}%
  \addtocounter{footnote}{-1}%
  \endgroup
}

\begin{document}
\author[Juan C. Arboleda Rivera]{Juan Camilo Arboleda Rivera}
\title{Error estándar}
%\subtitle{}
%\titlegraphic{\includegraphics[width=2cm]{udea}}
\institute[U. de A.]{Instituto de Biología\\
  Facultad de Ciencias Exactas y Naturales\\
  Universidad de Antioquia\\
  juan.arboleda2@udea.edu.co}
\date{v: \today}
%\subject{}

\begin{frame}[plain]
  \maketitle
\end{frame}

\begin{frame}
 \frametitle{Contenidos}
 \tableofcontents
\end{frame}

\section{Práctica Colab}

\begin{frame}
%[label=a]
{}{}
 \centering \LARGE
%
\href{https://colab.research.google.com/github/Nesper94/ciencia-de-datos/blob/main/notebooks/lec13-error-estándar.ipynb}
{\color{blue} Notebook Error Estándar}
\end{frame}


\section{¿Qué es el Error estándar?}

\begin{frame}
%[label=a]
{}{}
 \centering \Large

 \textbf{Error estándar}\\~\\

 El \textbf{error estándar} (SE por Standard Error) es la desviación
 estándar de la distribución muestral de un estadístico muestral, por
 ejemplo la media\\~\\

 El error estándar de la media sería la desviación estándar que tienen las
 medias muestrales tomadas de una población
\end{frame}

\begin{frame}
%[label=a]
{Cálculo del Error estándar de la media}{}
 \centering \large

 \textbf{\large Valor exacto}

 $$ \sigma_{\bar{x}} = \dfrac{\sigma}{\sqrt{n}}, $$

 donde $ \sigma $ es la desviación estándar \textbf{poblacional}, y $ n $ el
 tamaño de la muestra.\\~\\

\end{frame}

\begin{frame}
%[label=a]
{Cálculo del Error estándar}{}
 \centering \large
  \textbf{\large Estimación}\\~\\

  Como normalmente desconocemos $ \sigma $, se puede estimar el error
  estándar así:

  $$ \sigma_{\bar{x}} \approx \dfrac{s}{\sqrt{n}}, $$

  donde $ s $ es la desviación estándar \textbf{muestral}.
\end{frame}

\section{Ejercicios}

\begin{frame}
%[label=a]
{Crecimiento de plántulas}{}
 \centering \large

 Un estudiante mide la altura de 25 plántulas de caoba (\textit{Swietenia
 macrophylla}). Encuentra que el promedio de las alturas es 40 cm y la
 desviación estándar es de 5 cm\\~\\

 Calcula el error estándar de la media
\end{frame}

\begin{frame}
%[label=a]
{Masa de peces}{}
 \centering \large

 En una investigación sobre peces de agua dulce se pesan 36 individuos,
 dando un valor promedio de 90 g. La desviación estándar de la masa es de 12
 g\\~\\

 Calcula el error estándar de la media
\end{frame}

\begin{frame}
%[label=a]
{Frecuencia cardíaca en animales}{}
 \centering \large

 Un grupo de estudiantes registra la frecuencia cardíaca de ocho individuos
 de una especie, en condiciones experimentales controladas. Obtienen una
 media de 120 latidos/min y una desviación estándar de 8 latidos/min.\\~\\

 Responde:

 \begin{enumerate}

 \item ¿Cuál es el error estándar?

 \item ¿Qué representa este valor en el estudio?

 \item Si se aumenta la muestra a 32 individuos y la desviación estándar
 permanece en 8 latidos/min, ¿cuál sería el nuevo error estándar?
 \end{enumerate}

\end{frame}

\end{document}
